\section{Limitations and broader relevance}

Two biomedical datasets do not establish a universal law of conditional adaptation. RxRx1 and
Camelyon17 were chosen because both expose real acquisition environments while differing sharply in
modality, number of environments, and label structure. Environment metadata is available during
training for two pressure conditions and for analysis, which may not hold in every deployment.
Adversarial invariance can also remove task-relevant variation when environment and label are
statistically coupled.

The dense-substrate gate may select different pretrained ViTs for microscopy and histopathology.
Consequently, architectural effects are identified within each dataset against its matched dense
control; cross-dataset agreement is evidence of replication across regimes, while differences
cannot be attributed to acquisition modality alone. We prefer this limitation to running a common
but floor-performing backbone that would make the MoE comparison uninterpretable.

The broader value is an auditable way to decide whether conditional capacity is appropriate for
heterogeneous scientific adaptation. The design separates added capacity, routing, and invariance;
the diagnostics reveal when experts follow nuisance; and negative results define when a simpler
dense adaptation is preferable.
